\documentclass[11pt]{article}
\usepackage[utf8]{inputenc}
\usepackage[T1]{fontenc}
\usepackage{lmodern}
\usepackage[margin=1in]{geometry}
\usepackage{amsmath,amssymb,mathtools}
\usepackage{graphicx}
\usepackage{grffile}
\usepackage{float}
\usepackage{hyperref}
\usepackage{parskip}
\usepackage{enumitem}
\usepackage{xurl}
\setlength{\parskip}{0.65em}
\setlength{\parindent}{0pt}
\begin{document}
\title{Daily Research Digest --- 2026-04-17}
\author{}
\date{}
\maketitle
\textbf{Topics:} galaxy clusters, galaxies, gravitational lensing, dark matter.
\section*{Synthesis}
\_Digest generation failed: timed out\_
\newpage
\section*{Paper Summaries and Figures}
\subsection*{1. Age bimodality in pseudo-bulges of barred spiral galaxies: Bar-driven evolution across cosmic time}
\textbf{Focus:} galaxies\\
\textbf{Authors:} Kavin Kumar N R, Saili Keshri, Sudhanshu Barway\\
\textbf{Published:} 2026-04-16T17:35:25+00:00\\
\textbf{PDF:} \href{https://arxiv.org/pdf/2604.15257v1}{https://arxiv.org/pdf/2604.15257v1}\\
\textbf{Summary:}
We investigate the stellar population properties of pseudo-bulges in barred galaxies drawn from the Sloan Digital Sky Survey (SDSS DR7) to assess how bars regulate central star formation and secular evolution. Our sample comprises barred spiral and barred lenticular (S0) galaxies with reliable spectroscopic indices obtained from multicomponent structural decompositions. Stellar ages and recent star formation are traced using the 4000 Å break strength ($D_{n}(4000)$) and the Balmer absorption index ($H\delta _{A}$), complemented by bulge, bar, and disc colours. Barred spirals show a clear bimodality in $D_{n}(4000)$, with peaks at $D_{n}(4000)\sim1.3$ and $\sim1.8$. Low-$D_{n}(4000)$ pseudo-bulges exhibit strong $H\delta _{A}$ absorption, blue colours, and high specific star-formation rates, indicating young, actively growing centres. High-$D_{n}(4000)$ systems instead show weak $H\delta _{A}$, red colours, and low sSFR, consistent with older, quenched pseudo-bulges. Barred S0s display an old-bulge-dominated distribution, suggesting that gas-poor barred spirals transition into S0s following disc-wide quenching. We also find elevated AGN incidence among old pseudo-bulges. These trends support a scenario in which bars funnel gas inward to build pseudo-bulges and later suppress central star formation by depleting or stabilising the inflow. IFU observations show that bars assemble cold nuclear discs that age and quench over time, while high-redshift imaging confirms that bars are already present at $z\sim4$, implying that this evolutionary cycle operates across cosmic time. The strong correspondence between stellar age, colour, and structure indicates that bar-driven secular evolution governs both the growth and quenching of central components, linking blue barred spirals to red S0 galaxies.
\newpage
\subsection*{2. Echoes of Global Cosmic Strings}
\textbf{Focus:} dark matter\\
\textbf{Authors:} Jeff A. Dror, Antonios Kyriazis\\
\textbf{Published:} 2026-04-16T17:16:56+00:00\\
\textbf{PDF:} \href{https://arxiv.org/pdf/2604.15241v1}{https://arxiv.org/pdf/2604.15241v1}\\
\textbf{Summary:}
If the Universe underwent a cosmic phase transition, it may have left behind a network of cosmic strings. When these strings arise from the breaking of a gauge symmetry, their decay produces a significant stochastic background of gravitational waves. In contrast, if they originate from the breaking of a global symmetry, their decay predominantly yields Nambu-Goldstone bosons, which can persist as dark matter or dark radiation. In this work, we assess the detectability of this particle spectrum using a range of cosmological probes. We employ semi-numerical methods to estimate the resulting energy density and compute the associated matter power spectrum. We then compare these predictions with observations of the cosmic microwave background, Lyman-$\alpha $ forest, large-scale structure surveys, and the UV luminosity function, thereby deriving constraints on the Nambu-Goldstone boson mass and the symmetry-breaking scale. Finally, we present projections for the sensitivity of upcoming cosmic microwave background missions.
\subsubsection*{Selected figures}
\begin{figure}[H]
\centering
\includegraphics[width=\linewidth,height=0.42\textheight,keepaspectratio]{/home/kf/research-assistant/data/cache/figures/http-arxiv.org-abs-2604.15241v1/figure\_1.png}
\end{figure}
\newpage
\subsection*{3. Tracing evolutionary pathways of bar-driven quenching in local Universe disc galaxies}
\textbf{Focus:} galaxies\\
\textbf{Authors:} D Renu, Smitha Subramanian, Koshy George\\
\textbf{Published:} 2026-04-16T16:39:08+00:00\\
\textbf{PDF:} \href{https://arxiv.org/pdf/2604.15208v1}{https://arxiv.org/pdf/2604.15208v1}\\
\textbf{Summary:}
Bars play an integral role in regulating star formation (SF) in spiral galaxies, from triggering central starbursts to driving quenching. The diverse SF morphologies observed in local barred galaxies reflect different evolutionary stages of the bar, motivating studies across these stages. Here we study 12 nearby barred galaxies (z=0.01-0.06) identified as centrally quenched galaxies (having extended star-forming discs but quenched inner regions) by leveraging the differences in SFRs between the MPA-JHU and GSWLC catalogues. However, they exhibit residual central emission in the SDSS 3" fibre spectral region. Emission line analysis shows that this emission originates from either ongoing SF or LINER-like activity, suggesting diverse central ionization mechanisms. Using spatially resolved UV-optical colour maps from SDSS (r-band) and GALEX (FUV and NUV band) imaging data, we find that discs are star-forming and bluer in colour (NUV-r < 4 mag) while the bulge and bar regions are systematically redder (NUV-r > 4 mag) and dominated by older stellar populations. The NUV-r radial colour profiles show a clear transition from red to blue colours at the bar end with a corresponding median stellar age of \textasciitilde{} 1 Gyr. Compared to fully centrally quenched barred galaxies from our earlier work which lack SDSS fibre emission, these galaxies remain systematically bluer at similar radii, despite showing NUV-r > 4 mag inside the bar suggesting an intermediate stage in bar-driven quenching. We also estimate black hole masses associated with kinetic-mode AGN feedback and find them below the threshold (logM\_BH < 8.0). Adding this with the presence of pseudo bulges, our results support bars as the primary drivers of quenching, with these galaxies representing an evolutionary phase just before their inner regions are completely quenched.
\subsubsection*{Selected figures}
\begin{figure}[H]
\centering
\includegraphics[width=\linewidth,height=0.42\textheight,keepaspectratio]{/home/kf/research-assistant/data/cache/figures/http-arxiv.org-abs-2604.15208v1/figure\_1.png}
\end{figure}
\newpage
\subsection*{4. Understanding the regulation of star formation within TNG100 galaxies on kpc-scales using machine learning I: Global versus local}
\textbf{Focus:} galaxies\\
\textbf{Authors:} Bryanne McDonough, Sathvika S. Iyengar, Ansa Brew-Smith, Asa F. L. Bluck, Joanna Piotrowska\\
\textbf{Published:} 2026-04-16T16:10:55+00:00\\
\textbf{PDF:} \href{https://arxiv.org/pdf/2604.15182v1}{https://arxiv.org/pdf/2604.15182v1}\\
\textbf{Summary:}
We apply Random Forest and XGBoost machine learning algorithms to determine which galaxy properties most effectively predict star formation and quenching in simulated galaxies. Using spatially-resolved data from approximately 63,000 annular bins across 6,189 TNG100 galaxies, we train classification models to predict quenching states and regression models to predict star formation rate surface densities. Despite their different algorithmic approaches, both methods produce consistent feature importance rankings, with XGBoost distributing importance more evenly among correlated features. For central galaxies and high-mass satellites, black hole mass dominates quenching predictions, consistent with quenching via active galactic nuclei (AGN) feedback. Classification of low-mass satellites shows overwhelming importance for halo mass, indicating environmental quenching. Star formation predictions are dominated by local stellar mass surface density across all star-forming galaxy types, confirming that active star formation is a local process while quenching is driven by global properties.
\newpage
\subsection*{5. Cosmology of Inelastic Self-Interacting Dark Matter: Linear Evolution and Observational Constraints}
\textbf{Focus:} dark matter\\
\textbf{Authors:} Xin-Chen Duan, Yue-Lin Sming Tsai, Ziwei Wang\\
\textbf{Published:} 2026-04-16T13:33:16+00:00\\
\textbf{PDF:} \href{https://arxiv.org/pdf/2604.15006v1}{https://arxiv.org/pdf/2604.15006v1}\\
\textbf{Summary:}
We study the linear cosmological evolution of inelastic self-interacting dark matter in a two-component dark sector with a small mass splitting, assuming thermal initial conditions for the two species. We derive the coupled background and perturbation equations for inelastic conversion between the two species, considering both Power-law and Low-velocity saturation cross sections. Exothermic conversion injects kinetic energy into the light component, generating pressure support that suppresses small-scale structure and produces dark acoustic oscillations in the matter power spectrum. The resulting cutoff at scale $k > 1\,h\,\mathrm{Mpc}^{-1}$ depends on the normalization and velocity dependence of the cross section, the dark matter mass and the mass splitting. Using linear power spectra computed with a modified Boltzmann solver, we apply recast constraints from Lyman-$\alpha $ forest data and high-redshift UV luminosity functions, finding non-monotonic but closed exclusion regions driven by the competition between efficient conversion and rapid depletion of the heavy component. These results show that the internal thermodynamics of a secluded multi-component dark sector can leave observable imprints on structure formation, providing a complementary probe of dark matter beyond Standard Model interactions.
\newpage
\subsection*{6. Microscopic primordial black holes as macroscopic dark matter from large extra dimensions}
\textbf{Focus:} dark matter\\
\textbf{Authors:} Giuseppe Filiberto Vitale, Gaetano Lambiase, Tanmay Kumar Poddar, Luca Visinelli\\
\textbf{Published:} 2026-04-16T11:01:51+00:00\\
\textbf{PDF:} \href{https://arxiv.org/pdf/2604.14871v1}{https://arxiv.org/pdf/2604.14871v1}\\
\textbf{Summary:}
We study the coupled cosmological evolution of primordial black holes (PBHs) and radiation in the Arkani-Hamed-Dimopoulos-Dvali (ADD) framework with $n$ large extra dimensions and a fundamental gravity scale $M_\star$ at the TeV scale. For PBHs with horizon radius smaller than the compactification scale, the higher-dimensional geometry implies a larger horizon size at fixed mass and therefore a suppressed Hawking temperature. As a result, radiation accretion can overcome evaporation in the early Universe and drive a ``runaway'' phase of rapid mass growth. By numerically solving the coupled mass and energy-density evolution equations, we show that for $n \geq 2$ initially microscopic PBHs with initial mass $M_i \gtrsim 10^{12}\,$g can grow by many orders of magnitude and potentially reach macroscopic, even solar-mass, scales by matter-radiation equality. We determine the critical initial abundance $\beta _{\rm crit}$ required for PBHs to account for the observed dark matter density and find that extra dimensions dramatically lower this threshold, allowing viable scenarios with $\beta _{\rm crit}\sim 10^{-44}$. This identifies a previously unexplored region of parameter space in which the dark matter abundance is achieved through dynamical mass growth rather than large initial collapse fractions.
\subsubsection*{Selected figures}
\begin{figure}[H]
\centering
\includegraphics[width=\linewidth,height=0.42\textheight,keepaspectratio]{/home/kf/research-assistant/data/cache/figures/http-arxiv.org-abs-2604.14871v1/figure\_1.png}
\end{figure}
\newpage
\subsection*{7. Deeper analysis of Fermi-LAT unassociated 4FGL J2112.5-3043 for possible identification}
\textbf{Focus:} dark matter\\
\textbf{Authors:} Federica Giacchino, Cristina Fernández-Suárez, Miguel Á Sánchez-Conde, M. Ángeles Pérez-García, Stefano Ciprini, Dario Gasparrini\\
\textbf{Published:} 2026-04-16T08:55:20+00:00\\
\textbf{PDF:} \href{https://arxiv.org/pdf/2604.14794v1}{https://arxiv.org/pdf/2604.14794v1}\\
\textbf{Summary:}
In the 4FGL-DR4 point-source catalog of the Large Area Telescope (LAT) onboard NASA's Fermi Gamma-ray Observatory (Fermi-LAT), around a third of the sources are still unidentified (unIDs). In this work, we perform a detailed study of one of them, namely 4FGL J2112.5-3043. Only gamma-ray emission has been detected from this unidentified source, with no counterpart observed at any other wavelength as of today. Together with its high detection significance, this makes 4FGL J2112.5-3043 a particularly compelling target for further investigation. The results of our spectral and spatial analyses show that the source photon spectrum is better described with a subexponential cutoff power-law spectral model, with no significant flux variability over time, and a morphology consistent with being a point-like source. We investigate and discuss the characterized emission within the context of both conventional and exotic astrophysics, namely a pulsar origin or potential dark matter (DM) annihilations in a nearby Galactic subhalo. Although our results are inconclusive and neither confirm a DM origin nor firmly establish an astrophysical nature, we find a spectral preference for the $b\bar{b}$ and $c\bar{c}$ DM annihilation channels over a pulsar origin, thus making this unID a particularly intriguing candidate for next multiwavelength observations.
\subsubsection*{Selected figures}
\begin{figure}[H]
\centering
\includegraphics[width=\linewidth,height=0.42\textheight,keepaspectratio]{/home/kf/research-assistant/data/cache/figures/http-arxiv.org-abs-2604.14794v1/figure\_1.png}
\end{figure}
\newpage
\subsection*{8. Exploring non-equilibrium effects in sequential freeze-in}
\textbf{Focus:} dark matter\\
\textbf{Authors:} Shiuli Chatterjee, Andrzej Hryczuk\\
\textbf{Published:} 2026-04-16T06:47:20+00:00\\
\textbf{PDF:} \href{https://arxiv.org/pdf/2604.14688v1}{https://arxiv.org/pdf/2604.14688v1}\\
\textbf{Summary:}
Freeze-in of multi-component dark sectors is governed not only by the interaction with the thermal plasma, but also by their internal dynamics. Full thermalisation within the dark sector is not guaranteed, raising the question of impact of departures from local thermal equilibrium onto the evolution and ultimately relic abundance and momentum distribution of dark matter. In this work we explore this question in a minimal two-scalar model, which can give rise to observable signatures in indirect detection and long-lived particle searches at forward physics experiments. Focusing on the phenomenologically viable regions, we analyse the impact of non-thermal evolution on the dark matter abundance, finding deviations of up to an order of magnitude between the full phase-space treatment and the traditional number-density approach. Our results highlight the importance of phase-space level computation for accurate freeze-in predictions and further motivate dedicated numerical tools for studying the evolution of multi-component dark sectors at the phase space level.
\newpage
\subsection*{9. Inflaton Regeneration via Scalar Couplings: Generic Models and the Higgs Portal}
\textbf{Focus:} dark matter\\
\textbf{Authors:} Kunio Kaneta, Tomo Takahashi, Natsumi Watanabe\\
\textbf{Published:} 2026-04-16T05:02:37+00:00\\
\textbf{PDF:} \href{https://arxiv.org/pdf/2604.14620v1}{https://arxiv.org/pdf/2604.14620v1}\\
\textbf{Summary:}
The standard cosmological paradigm assumes that the inflaton field becomes dynamically negligible during the post-reheating evolution of the Universe. We demonstrate that this assumption fails for a broad class of inflationary models where the potential behaves as a monomial form $V(\phi ) \propto \phi ^k$ (with $k \ge 4$) around the minimum. In such scenarios, the effective inflaton mass depends on the field amplitude and vanishes asymptotically as the Universe expands. This vanishing-mass mechanism renders the inflaton kinematically accessible to the thermal plasma long after reheating, facilitating the regeneration of inflaton quanta through 1-to-2 decays and 2-to-2 scatterings of bath particles. This mechanism is quite generic and the coupling responsible for reheating can be constrained if the inflaton is overproduced, while the inflaton quanta can constitute dark matter in specific scenarios. Furthermore, if reheating occurs via the Standard Model Higgs portal, the process can be further constrained by big bang nucleosynthesis, cosmic microwave background, and colliders such as the LHC. This mechanism provides a new framework for probing post-inflationary reheating.
\newpage
\subsection*{10. Gravitational Lensing Signatures of Hayward-like Black Holes}
\textbf{Focus:} gravitational lensing\\
\textbf{Authors:} Chen-Hung Hsiao, Limei Yuan, Yidun Wan\\
\textbf{Published:} 2026-04-16T00:38:35+00:00\\
\textbf{PDF:} \href{https://arxiv.org/pdf/2604.14505v1}{https://arxiv.org/pdf/2604.14505v1}\\
\textbf{Summary:}
We examine the gravitational lensing signatures of a Hayward-like regular black hole and its potential observational distinction from a Schwarzschild black hole. In the weak-field limit, the deflection angle includes a small positive correction proportional to $m \ell^2/b^3$, indicating slightly stronger light bending than in Schwarzschild, though the effect remains observationally negligible at large impact parameters. Current galaxy-scale Einstein-ring data, such as from ESO325-G004, cannot yet constrain the regular-core scale $\ell$. In the strong-deflection regime, for Sgr A* and M87*, the asymptotic position $θ_{\infty}$ is identical to Schwarzschild's. Nevertheless, $\ell$ modifies strong-lensing coefficients $\bar a, \bar b$, influencing angular separations s, relative flux ratio $r_\mathrm{mag}$, and time delays $\Delta T_{2,1}$. Our predicted values for these observables remain consistent with current data, suggesting that future high-precision measurements of strong-field lensing may distinguish Hayward-like from Schwarzschild black holes.
\subsubsection*{Selected figures}
\begin{figure}[H]
\centering
\includegraphics[width=\linewidth,height=0.42\textheight,keepaspectratio]{/home/kf/research-assistant/data/cache/figures/http-arxiv.org-abs-2604.14505v1/figure\_1.png}
\end{figure}
\end{document}
